\section{Limitations}

The breadth of the framework does not make value functions interchangeable
across domains. Users remain responsible for model structure, causal
assumptions, input provenance, outcome valuation, thresholds, distributional
choices, and interpretation. A monetary net-benefit model, an equity-weighted
social objective, and an engineering loss function can share computational
structure while representing different ethical and empirical claims.

Regression-based EVPPI and EVSI estimates inherit modelling choices and Monte
Carlo error. Convergence, effective sample size, predictive adequacy,
extrapolation, and sensitivity to method settings provide useful context for a
scalar result. Population VOI additionally depends on
eligible population, time horizon, discounting, implementation, and research
cost assumptions; \texttt{voiage} validates representations and numerical
preconditions but cannot determine whether those assumptions are substantively
appropriate.

Newer methods are individually marked as developing or experimental. Their
repeatable examples show that the stated calculations and software behaviour
have been implemented. They do not establish statistical validity, accuracy in
real applications, superiority to other methods, suitability in another
setting, or usefulness for policy. In particular,
perspective, data-quality, transportability, dynamic-option, and several
evidence-system methods need applied comparisons and domain-specific validation.
The Python package is broader than the current R and Julia packages. Both
provide native EVPI and ENBS in version 2.2.0, but neither exposes Python's
full decision-record interface. R's optional EVPPI and EVSI paths require Python.

Automated checks detect some implementation errors but do not replace
independent statistical validation, expert review, sensitivity analysis, or
replication with relevant domain data. Optional accelerators and distributed
or hardware research are not evidence of faster production analysis.
